\section{Source Language}
\label{sec:sourcelanguage}

One source file represents a project and its applications.
Table~\ref{tab:source_language} summarises the language.

\begin{table}[!htbp]
\centering
\small
\begin{tabular}{>{\raggedright\arraybackslash}p{0.24\linewidth}>{\raggedright\arraybackslash}p{0.68\linewidth}}
\hline
\textbf{Concern} & \textbf{Model content} \\
\hline
Ownership and release &
  Who owns the project and which processes are released together as one application. \\
Processes &
  Which containerised processes run and whether they are one-time or long-running. \\
Placement and scale &
  A process's CPU, memory, allowed machines and number of running instances. \\
Exposure and routing &
  An application's web address, who may reach it and which process handles each URL path. \\
Connections &
  The ports a process offers and the applications it must reach. \\
Health and cutover &
  Checks for process health and readiness, and how a new version replaces an old one. \\
Observability &
  Where the platform collects metrics and how urgently a failure raises an alert. \\
Storage and writes &
  Where a process stores data, how it is recovered and which folders it may access. \\
Secrets &
  Which vault secrets a process may read and how it receives them. \\
\hline
\end{tabular}
\caption{Concerns Defined by the Source Language}
\label{tab:source_language}
\end{table}

The model describes application needs rather than Kubernetes kinds, proxy
fields or machine names. The model-to-model transformation adds the
platform-specific details before the final files are generated.
